\subsection{9. Motion Tomography via Occupation Kernels}
\textbf{OSTI:} \texttt{1842593}\quad\textbf{Authors:} Christensen et al. (2021)\quad\textbf{Domain:} Control Theory / Kernel Methods\\
\scorebadge{Coverage}{8}\quad\scorebadge{Agreement}{8}\quad\textbf{Total:} 16/20\par\smallskip
\noindent\textbf{Coverage rationale.} Full Algorithm 1 implementation (iterative flow-field reconstruction via occupation-kernel Gram-matrix inversion). Three synthetic flow fields (Gaussian mixture, linear/spiral, constant) tested. Parameter sensitivity sweeps over $\mu$ (kernel width), $\lambda$ (regularization), N (number of trajectories). 26/26 tests pass covering kernel properties, Gram-matrix structure, and reconstruction convergence. Main gap: Experiment 2 (real-world Gliderpalooza dataset from Chang et al. [5]) not replicated --- no data access. No comparison with alternative RKHS-based or deep-learning flow reconstruction methods.\par
\noindent\textbf{Agreement rationale.} Algorithm converges with 5-order-of-magnitude displacement reduction over 10 iterations (0.265 $\rightarrow$ 5$\times$10). Mean relative error 0.0155 (paper: 0.0256) --- ours slightly better, attributed to different random trajectory seeds. Convergence ordering Constant \textless{} Linear \textless{} Gaussian confirmed. U-shaped $\mu$ sensitivity with optimum $\approx$ 0.5 confirmed. N sweep shows monotonic improvement with more trajectories. $\lambda$ sensitivity monotonically increasing (paper's was U-shaped, likely because of different Gram-matrix conditioning).\par\smallskip
\noindent\textbf{What reproduced:}
\begin{itemize}[leftmargin=1.4em,itemsep=0.2em,topsep=0.2em]
\item Algorithm 1 iterative flow reconstruction via occupation-kernel RKHS
\item Gaussian RBF and exponential-dot-product kernels
\item Occupation-kernel evaluation via Simpson's rule
\item Gram-matrix computation via 2D Simpson's rule
\item All three synthetic flow fields from the paper
\item Parameter sweeps over kernel width $\mu$, regularization $\lambda$, trajectory count N
\item Convergence ordering and rates for the three fields
\item Error metric (mean relative error) within paper's range
\end{itemize}\noindent\textbf{What missing:}
\begin{itemize}[leftmargin=1.4em,itemsep=0.2em,topsep=0.2em]
\item Experiment 2: Real-world Gliderpalooza dataset from Chang et al. [5]
\item Comparison with alternative flow-reconstruction methods (sparse regression, GP, NN)
\item Theoretical convergence-rate verification against Theorems 2.3, 3.1
\item Robustness to observation noise in endpoint measurements
\item Non-Gaussian kernel families beyond RBF and exponential-dot-product
\item Scaling analysis: wall-clock vs N (our O(N$^2$$\cdot$n$^2$) cost model)
\item Time-varying flow fields F(x,t)
\end{itemize}\noindent\textbf{Follow-on research questions:}
\begin{enumerate}[leftmargin=1.6em,itemsep=0.2em,topsep=0.2em,label=Q\arabic*.]
\item Download a public oceanographic glider dataset (e.g., IOOS Glider DAC) and apply Algorithm 1 to real endpoint displacements --- does kernel width $\mu$=0.5 remain optimal on physical-scale flows, and what is the reconstruction error vs HYCOM model ground truth
\item Add Gaussian observation noise $\sigma$ to endpoint displacements and characterize algorithm breakdown --- at what $\sigma$/displacement ratio does the reconstruction error exceed 20\%, and does Tikhonov regularization extend the tolerable noise range
\item Replace the occupation-kernel RKHS with a divergence-free neural-network parameterization --- does the learned approach beat the kernel method on the Gaussian-mixture flow at N=25, and how does it scale to 4D (spatial+time) flows
\item Extend Algorithm 1 to time-varying flows F(x,t) via space-time kernels K((x,t),(y,s)) --- does it recover a known oscillating-vortex test case, and what is the effective temporal resolution as a function of trajectory density
\item For the Gaussian-mixture flow, characterize identifiability: with fixed trajectory endpoints, what subspace of F is unrecoverable, and does this null space correspond to a closed-form RKHS orthogonal complement
\end{enumerate}

